\section{Introdução}
\label{sec:introducao}

O teste de software é uma das principais atividades para garantir a qualidade e a confiabilidade de um
sistema \cite{ahamed2010}. Dentro dele, o teste não funcional, que verifica desempenho, segurança,
confiabilidade, usabilidade e outras características de qualidade, ainda recebe pouco espaço nos cursos de
graduação, apesar da sua relevância na indústria \cite{garousi2020}. Diferentemente do teste funcional, que
verifica se o sistema faz o que foi especificado, o teste não funcional avalia \emph{como} o sistema se comporta
sob diferentes condições, o que o torna mais abstrato e mais difícil de ensinar \cite{paschoal2018}. Aulas
expositivas e exercícios isolados costumam ser insuficientes para que os estudantes liguem as definições dos
livros a decisões práticas \cite{valle2015,vanhanen2018}.

Jogos educacionais vêm sendo usados para enfrentar esse problema no ensino de teste de software
\cite{costa2024debug,gomes2021gamut,beppe2018greatest,jesus2020}. Nesse contexto, a primeira versão do
\emph{Critical Failure} foi criada na disciplina de Teste de Software do Centro de Informática da UFPB
\cite{oliveira2025critical}. Nela, os jogadores montam um ``sistema'' com cartas de testes não funcionais, testam
os sistemas dos colegas e podem blefar ao se defender. Dezessete estudantes jogaram e avaliaram o jogo de forma
positiva, mas a aplicação também revelou um problema de projeto: em quatro das seis partidas observadas, o
vencedor terminou com pontuação muito acima dos demais. Quem saía na frente acumulava vantagem, e os
jogadores tendiam a atacar os mais fracos, e não o líder \cite{oliveira2025critical}. O próprio artigo sugeriu
investigar se esse comportamento vinha da inexperiência dos jogadores ou de uma falha da dinâmica do jogo.

Corrigir o balanceamento de um jogo pede muitas partidas. Com pessoas, cada partida leva quase uma hora, e
algumas dezenas de partidas não bastam para separar o efeito de uma regra do acaso, do estilo dos jogadores ou
da ordem à mesa. Este trabalho trata o jogo como um \emph{sistema sob teste}: o balanceamento é escrito como um
conjunto de requisitos não funcionais com critério de aceite, e um simulador joga milhares de partidas entre
bots para verificar esses requisitos antes de levar as regras de volta à sala de aula. Assim, o trabalho aplica
teste não funcional ao próprio jogo que ensina teste não funcional.

\subsection{Objetivo e perguntas de pesquisa}

O objetivo deste trabalho é redesenhar as regras do Critical Failure e verificar, por simulação, se a nova versão
é balanceada, comparando estilos de jogo diferentes para identificar o melhor estilo e as fraquezas das regras.
Duas perguntas de pesquisa orientam o trabalho:

\begin{itemize}
  \item \textbf{QP1:} as novas regras reduzem o desequilíbrio observado na primeira versão?
  \item \textbf{QP2:} como os estudantes avaliam a experiência com a segunda versão?
\end{itemize}

\subsection{Contribuições}

As contribuições deste trabalho são:
\begin{enumerate}
  \item a segunda versão das regras do Critical Failure (v2.2.2), alinhada à ISO/IEC 25010:2023
    \cite{iso25010}, com cada correção e cada interpretação registradas;
  \item um simulador de partidas determinístico e auditável, em que toda partida pode ser refeita e cada evento
    cita a regra que o justificou;
  \item um conjunto de 13 requisitos de balanceamento com zonas de aceite fixadas antes da medição, e a
    verificação deles em 971 mil partidas;
  \item a comparação de seis estilos de jogo e de perfis de blefe, de variantes, da retirada de cada mudança
    de regra (ablação) e da sensibilidade dos parâmetros;
  \item a avaliação de um playtest com estudantes, ligada aos resultados do simulador.
\end{enumerate}

\subsection{Organização do texto}

A Seção~\ref{sec:relacionados} apresenta os trabalhos relacionados. A Seção~\ref{sec:metodologia} descreve a
metodologia, baseada em \emph{Design Science Research}. A Seção~\ref{sec:implementacao} apresenta as regras da
segunda versão e o simulador. A Seção~\ref{sec:resultados} traz os resultados, as lições aprendidas e as
limitações. A Seção~\ref{sec:conclusoes} conclui o trabalho e aponta trabalhos futuros.
